\section{Desain Sistem}
\label{sec:design}

\subsection{Gambaran Arsitektur}
Gambar~\ref{fig:arch} mengilustrasikan arsitektur HalalChain melintasi penyimpanan dokumen off-chain dan pengaitan status on-chain.
Produsen mengunggah dokumen pendukung ke IPFS melalui integrasi Pinata sisi server, memperoleh CID, dan memanggil \texttt{registerBatch} pada Base L2 dengan CID tersebut.
Auditor mengambil dokumen melalui gateway IPFS, mencatat keputusan on-chain, dan secara opsional melampirkan CID laporan audit.
Konsumen mengkueri status batch secara read-only melalui \texttt{getBatch} via endpoint JSON-RPC yang disematkan dalam URL web; kode QR mengodekan URL ini sehingga pemindaian tidak memerlukan perangkat lunak dompet.

\begin{figure}[!t]
  \centering
  \resizebox{\columnwidth}{!}{%
  \begin{tikzpicture}[
    font=\footnotesize,
    >=Stealth,
    actor/.style={
      draw,
      rounded corners=2pt,
      line width=0.4pt,
      fill=gray!6,
      minimum width=2.05cm,
      minimum height=0.9cm,
      align=center,
      inner sep=4pt
    },
    core/.style={
      draw,
      rounded corners=2pt,
      line width=0.5pt,
      fill=gray!14,
      minimum width=2.35cm,
      minimum height=0.95cm,
      align=center,
      inner sep=4pt
    },
    store/.style={
      draw,
      rounded corners=2pt,
      line width=0.4pt,
      fill=gray!10,
      minimum width=2.05cm,
      minimum height=0.9cm,
      align=center,
      inner sep=4pt
    },
    lbl/.style={
      font=\scriptsize,
      fill=white,
      inner sep=1.5pt,
      text=black
    }
  ]
    \node[actor] (producer) {Produsen\\UMKM};
    \node[store, right=1.25cm of producer] (ipfs) {IPFS\\Pinata};
    \node[core] at ($(producer)!0.5!(ipfs) + (0,-1.45)$) (chain) {Base L2\\HalalChain.sol};

    \node[actor, below left=1.05cm and 0.55cm of chain] (consumer) {Konsumen\\verifikasi QR};
    \node[actor, below right=1.05cm and 0.55cm of chain] (auditor) {Auditor};

    \draw[->, thick] (producer.east) -- node[lbl, above]{unggah dok.} (ipfs.west);
    \draw[->, thick] (producer.south) -- ++(0,-0.55) -|
      node[lbl, pos=0.2, below]{registerBatch} (chain.west);

    \draw[->, thick, dashed] (ipfs.south) -- node[lbl, right]{simpan CID} (chain.north);

    \draw[->, thick] (chain.south west) -- node[lbl, left, pos=0.35]{baca status} (consumer.north);
    \draw[->, thick] (auditor.north) -- node[lbl, right, pos=0.35]{verifikasi / tolak} (chain.south east);

  \end{tikzpicture}%
  }
  \caption{Arsitektur HalalChain: dokumen pada IPFS, jangkar dan status pada Base L2.}
  \label{fig:arch}
\end{figure}


Desain memisahkan \emph{bukti} (hosting besar dan dapat berubah pada IPFS) dari \emph{komitmen} (kecil, direplikasi, permanen pada L2).
Tabel~\ref{tab:datasplit} merangkum pemisahan tersebut.

\begin{table}[t]
  \caption{Penempatan data on-chain vs.\ off-chain.}
  \label{tab:datasplit}
  \centering
  \footnotesize
  \begin{tabularx}{\columnwidth}{@{}>{\raggedright\arraybackslash}p{0.36\columnwidth}>{\raggedright\arraybackslash}p{0.18\columnwidth}X@{}}
    \toprule
    \textbf{Data} & \textbf{Lokasi} & \textbf{Alasan} \\
    \midrule
    Status batch, cap waktu & On-chain & FSM tahan manipulasi \\
    Alamat produsen/auditor & On-chain & Akuntabilitas \\
    CID IPFS & On-chain & Jangkar berbasis alamat konten \\
    Alasan penolakan (string) & On-chain & Jejak audit permanen \\
    PDF, gambar, laporan lab & IPFS & Biaya dan ukuran \\
    Berkas pendukung audit & IPFS & Sama seperti dokumen produsen \\
    \bottomrule
  \end{tabularx}
\end{table}

\subsection{Siklus Hidup Batch}
Status batch mengikuti mesin status terbatas (FSM) dengan lima nilai enumerasi: \texttt{Unknown}, \texttt{Pending}, \texttt{Verified}, \texttt{Rejected}, dan \texttt{Revoked}.
Tabel~\ref{tab:fsm} mencantumkan transisi yang diizinkan; semua transisi lain gagal pada lapisan kontrak.

\begin{table}[t]
  \caption{Transisi status batch yang ditegakkan oleh \texttt{HalalChain.sol}.}
  \label{tab:fsm}
  \centering
  \footnotesize
  \setlength{\tabcolsep}{3pt}
  \begin{tabularx}{\columnwidth}{@{}>{\raggedright\arraybackslash}p{0.14\columnwidth}>{\raggedright\arraybackslash}p{0.20\columnwidth}>{\raggedright\arraybackslash}X@{}}
    \toprule
    \textbf{Dari} & \textbf{Ke} & \textbf{Pemicu} \\
    \midrule
    --- & Pending & \makecell[tl]{\texttt{registerBatch},\\\texttt{registerRevision}} \\
    Pending & Verified & \makecell[tl]{\texttt{verifyBatch}\\\textit{(auditor)}} \\
    Pending & Rejected & \makecell[tl]{\texttt{rejectBatch}\\\textit{(auditor)}} \\
    Verified & Revoked & \makecell[tl]{\texttt{revokeBatch}\\\textit{(auditor)}} \\
    Rejected & Pending (id baru) & \texttt{registerRevision} \\
    Rejected & Verified & \emph{Dilarang} (gagal) \\
    \bottomrule
  \end{tabularx}
\end{table}

Batch yang ditolak tidak dapat diverifikasi ulang pada tempatnya; produsen mengajukan \emph{revisi} sebagai batch baru yang ditautkan melalui \texttt{parentBatchId}, mempertahankan alasan penolakan asli dan metadata audit untuk pemeriksaan.
Pencabutan hanya berlaku untuk batch yang sebelumnya Verified---misalnya, ketika pengawasan pasca-pasar menemukan kecurangan---dan mencatat string alasan on-chain.

\subsection{Spesifikasi Kontrak Pintar}
Kontrak \texttt{HalalChain.sol} (Solidity ${\hat{}}0.8.20$, OpenZeppelin \texttt{AccessControl}~\cite{openzeppelin2024access}) mendefinisikan konstanta peran \texttt{PRODUCER\_ROLE}, \texttt{AUDITOR\_ROLE}, dan \texttt{DEFAULT\_ADMIN\_ROLE}.
Kesalahan kustom (\texttt{InvalidBatch}, \texttt{InvalidStatus}, \texttt{NotOriginalProducer}, \texttt{ParentNotRejected}) menggantikan string revert untuk menghemat gas dan meningkatkan kejelasan integrator.

\textbf{Field struct Batch:}
\texttt{producer}, \texttt{productName}, \texttt{ipfsCid}, \texttt{createdAt}, \texttt{status}, \texttt{auditor}, \texttt{auditedAt}, \texttt{auditIpfsCid}, \texttt{rejectReason}, \texttt{parentBatchId}.

\textbf{Fungsi utama:}
\begin{itemize}
  \item \texttt{registerBatch(productName, ipfsCid)} --- membuat batch dengan \texttt{parentBatchId = 0}, status Pending; memancarkan \texttt{BatchRegistered}.
  \item \texttt{registerRevision(parentBatchId, ipfsCid)} --- mensyaratkan induk Rejected dan \texttt{msg.sender == parent.producer}; mewarisi \texttt{productName}; memancarkan peristiwa dengan tautan induk.
  \item \texttt{verifyBatch(batchId, auditIpfsCid)} --- Pending $\rightarrow$ Verified; mengatur field auditor; memancarkan \texttt{BatchVerified}.
  \item \texttt{rejectBatch(batchId, reason, auditIpfsCid)} --- Pending $\rightarrow$ Rejected; menyimpan alasan; memancarkan \texttt{BatchRejected}.
  \item \texttt{revokeBatch(batchId, reason)} --- Verified $\rightarrow$ Revoked; memancarkan \texttt{BatchRevoked}.
  \item \texttt{getBatch(batchId)} --- view publik mengembalikan struct lengkap; digunakan halaman konsumen tanpa biaya transaksi.
\end{itemize}

Peristiwa diindeks pada \texttt{batchId} dan alamat peserta jika berlaku, memungkinkan block explorer dan pengindeks subgraph merekonstruksi linimasa audit tanpa API proprietari.

\subsection{Penyimpanan Off-Chain dan Integritas}
Dokumen dipin pada IPFS melalui Pinata; hanya CID yang disimpan on-chain~\cite{benet2014ipfs,christidis2016blockchain}.
Integritas mengikuti penandaan alamat konten: jika gateway mengembalikan byte yang hash-nya tidak cocok dengan CID yang dikaitkan, klien harus memperlakukan dokumen sebagai tidak terpercaya.
Ketersediaan bergantung pada kebijakan pinning---CID yang tidak dipin dapat menjadi tidak dapat dijangkau meskipun referensi on-chain tetap ada---keterbatasan yang dianalisis di Bagian~\ref{sec:discussion}.
API unggah tidak pernah mengekspos kredensial JWT Pinata ke browser; hanya rute Next.js sisi server yang menyimpan rahasia.

\subsection{Jalur Verifikasi Konsumen}
Konsumen mengakses \texttt{/verify/\{batchId\}} pada frontend yang di-deploy.
Halaman memanggil \texttt{getBatch} melalui endpoint RPC publik (klien baca Viem/Wagmi), merender UI spesifik status (Verified / Pending / Rejected / Revoked), menampilkan tautan IPFS untuk dokumen produsen dan audit, serta menampilkan \texttt{parentBatchId} jika ada.
Tidak ada materi kunci privat yang dimuat pada rute konsumen, memenuhi persyaratan tanpa dompet pada Tabel~\ref{tab:related}.
